\begin{frame}{10.3 마무리: ``절대 보상''이 아니라 ``평균 대비
향상''}
  \begin{block}{핵심 정리}
    REINFORCE의 신호 $G_t$ 에는 \textbf{``운''}이 섞여 있다.
    \textbf{베이스라인} $b(s_t)$ --- \textbf{행동에 무관한 값} ---
    을 빼도 그래디언트의 \textbf{기댓값은 그대로}(행동의 확률
    합이 1에서), \textbf{분산만 줄어듦}.
    \vskip0.2em
    분산을 최소화하는 베이스라인은 $V(s_t)$ 자체 ---
    따라서 \textbf{가치함수를 \kb{Critic}}이라 부르고,
    $G_t - V(s_t)$ 를 \textbf{\kb{어드밴티지}}라 부르며,
    Actor는 이 어드밴티지로, Critic은 \textbf{TD(0)}으로,
    \textbf{동시에} 학습.
  \end{block}
  \vskip0.4em
  \kq{``절대 보상''이 아니라 \textbf{``평균 대비 향상''}으로
  배우는 이 한 가지 아이디어가 Week 11의 PPO로 이어진다.}
\end{frame}
